\section{Квадратичные формы. Полярные билинейные формы. Метод Лагранжа приведения квадратичной формы к каноническому и нормальному виду. Метод Якоби}

\begin{definition}
Функция $Q: V \to F$ называется \textbf{квадратичной формой}, если существует симметричная билинейная форма $f$ такая, что для любого $x \in V$ выполняется:
\[
Q(x) = f(x, x).
\]
Форма $f$ называется \textbf{полярной} (или поляризующей) билинейной формой для $Q$.
\end{definition}

\begin{theorem}[О восстановлении полярной формы]
Симметричная билинейная форма $f$, порождающая квадратичную форму $Q$, восстанавливается по $Q$ единственным образом с помощью формулы:
\[
f(x, y) = \frac{1}{2} \big( Q(x+y) - Q(x) - Q(y) \big).
\]
\end{theorem}

\begin{proof}
Используя тождество поляризации из билета 34:
\[
Q(x+y) = f(x+y, x+y) = f(x, x) + 2f(x, y) + f(y, y) = Q(x) + 2f(x, y) + Q(y).
\]
Выражая отсюда $f(x, y)$, получаем требуемую формулу. Единственность следует из того, что симметричная билинейная форма однозначно определяется своими значениями на диагонали.
\hfill$\blacksquare$
\end{proof}

\begin{definition}
Квадратичная форма называется \textbf{канонической}, если в некотором базисе она имеет вид суммы квадратов координат:
\[
Q(x) = \alpha_1 \xi_1^2 + \alpha_2 \xi_2^2 + \dots + \alpha_n \xi_n^2, \quad \alpha_i \in F.
\]
Если все $\alpha_i \in \{1, -1, 0\}$, форма называется \textbf{нормальной}.
\end{definition}

\begin{theorem}[Метод Лагранжа]
Любую квадратичную форму над полем $F$ ($\operatorname{char} F \neq 2$) можно привести к каноническому виду с помощью невырожденного линейного преобразования координат.
\end{theorem}

\begin{proof}
Доказательство проводится индукцией по размерности пространства $n$.
\begin{enumerate}
    \item \textbf{База индукции:} При $n=1$ форма уже имеет вид $Q(x) = a_{11} \xi_1^2$, что является каноническим.
    \item \textbf{Индукционный переход:} Пусть утверждение верно для размерности $n-1$. Рассмотрим форму $Q(x) = \sum_{i,j=1}^n a_{ij} \xi_i \xi_j$.
    
    \textit{Случай 1:} Существует $a_{kk} \neq 0$. Без ограничения общности (перестановкой переменных) считаем $a_{11} \neq 0$. Выделим все слагаемые, содержащие $\xi_1$:
    \[
    Q(x) = a_{11} \xi_1^2 + 2 \sum_{j=2}^n a_{1j} \xi_1 \xi_j + R(\xi_2, \dots, \xi_n),
    \]
    где $R$ --- квадратичная форма от $n-1$ переменной. Дополним выражение до полного квадрата:
    \[
    Q(x) = a_{11} \left( \xi_1 + \sum_{j=2}^n \frac{a_{1j}}{a_{11}} \xi_j \right)^2 - a_{11} \left( \sum_{j=2}^n \frac{a_{1j}}{a_{11}} \xi_j \right)^2 + R(\xi_2, \dots, \xi_n).
    \]
    Введем новую переменную $\eta_1 = \xi_1 + \sum_{j=2}^n \frac{a_{1j}}{a_{11}} \xi_j$, а остальные $\eta_j = \xi_j$ ($j \ge 2$). Тогда $Q(x) = a_{11} \eta_1^2 + Q_1(\eta_2, \dots, \eta_n)$. По предположению индукции, $Q_1$ приводится к каноническому виду.
    
    \textit{Случай 2:} Все $a_{ii} = 0$, но существует $a_{ij} \neq 0$ при $i \neq j$. Без ограничения общности пусть $a_{12} \neq 0$. Сделаем невырожденную подстановку:
    \[
    \xi_1 = \eta_1 + \eta_2, \quad \xi_2 = \eta_1 - \eta_2, \quad \xi_k = \eta_k \ (k \ge 3).
    \]
    Тогда слагаемое $2a_{12} \xi_1 \xi_2$ превратится в $2a_{12} (\eta_1^2 - \eta_2^2)$. Поскольку $a_{12} \neq 0$, мы получили квадраты с ненулевыми коэффициентами при $\eta_1^2$ и $\eta_2^2$, и задача сводится к Случаю 1.
\end{enumerate}
В обоих случаях преобразование координат невырождено, так как оно задается треугольной или блочной матрицей с ненулевым определителем.
\hfill$\blacksquare$
\end{proof}

\begin{theorem}[Метод Якоби]
Пусть $A$ --- матрица квадратичной формы $Q$ в некотором базисе, и все ее главные угловые миноры $\Delta_k$ ($k = 1, \dots, n$) отличны от нуля ($\Delta_0 = 1$ по определению). Тогда существует невырожденное линейное преобразование, приводящее форму к каноническому виду:
\[
Q(x) = \frac{\Delta_1}{\Delta_0} \eta_1^2 + \frac{\Delta_2}{\Delta_1} \eta_2^2 + \dots + \frac{\Delta_n}{\Delta_{n-1}} \eta_n^2.
\]
\end{theorem}

\begin{remark}
Доказательство этого факта опускается (согласно условию), но оно опирается на существование треугольного разложения матрицы $A = L D L^T$ (разложение Холецкого для симметричных матриц с ненулевыми главными минорами), где $D$ --- диагональная матрица с элементами $d_k = \Delta_k / \Delta_{k-1}$.
\end{remark}
